% Exemplo ILUSTRATIVO de frame de comparação com narrativa — copiar a estrutura, não o
% conteúdo (números e decisão abaixo são fictícios, só para mostrar o formato).
% Princípio: cada frame conta uma HISTÓRIA de 3 batidas —
% (1) o que medimos, (2) o que isso significou, (3) o que decidimos e por quê —
% nunca só uma tabela solta sem a decisão que ela gerou.

\begin{frame}{Escolha de banco de dados: Postgres vs. MongoDB}
  \begin{numerobloco}{O que medimos (benchmark de carga, 10k escritas/leituras)}
    \begin{center}
    \begin{tabular}{lccc}
      \toprule
      \textbf{Métrica} & \textbf{Postgres} & \textbf{MongoDB} & \textbf{Postgres + índice extra} \\
      \midrule
      Latência p50 (ms) & 12   & 9    & 8    \\
      Latência p99 (ms) & 45   & 38   & \alert{22}    \\
      Consistência forte & Sim  & Não  & Sim  \\
      \bottomrule
    \end{tabular}
    \end{center}
  \end{numerobloco}

  \pause
  \textbf{O que isso significou:} o ganho de latência do MongoDB vinha da ausência de
  consistência forte, não de uma vantagem estrutural do banco — replicando o mesmo índice
  extra no Postgres, a diferença de p99 \alert{some} e ainda mantemos consistência forte.

  \pause
  \begin{decisaobloco}{O que decidimos}
    Ficamos com Postgres + índice extra: a latência empata com a melhor opção, e não abrimos
    mão de consistência forte, que é um requisito do domínio (transações financeiras).
  \end{decisaobloco}
\end{frame}
